% =============================================================================
%  entete.tex — Personnalisation LaTeX commune à tous les PDF de cours
%  (inclus automatiquement par Quarto via `include-in-header`)
% =============================================================================

% --- Polices --------------------------------------------------------------------
% Atkinson Hyperlegible (texte) et DejaVu Sans Mono (code). Les fichiers sont
% fournis dans latex/polices/ : LuaLaTeX les trouve si ce dossier est indiqué
% dans la variable d'environnement OSFONTDIR (fait automatiquement par le
% workflow GitHub) ou si les polices sont installées sur l'ordinateur.
% Si une police est introuvable, on se replie sur une autre : la compilation
% n'échoue jamais pour une question de police.
\IfFontExistsTF{Atkinson Hyperlegible}{%
  \setmainfont{Atkinson Hyperlegible}%
  \setsansfont{Atkinson Hyperlegible}%   titres, en-têtes
}{%
  \IfFontExistsTF{DejaVu Sans}{%
    \setmainfont{DejaVu Sans}\setsansfont{DejaVu Sans}}{%
    \renewcommand{\familydefault}{\sfdefault}}%
}
\IfFontExistsTF{DejaVu Sans Mono}{%
  \setmonofont[Scale=0.9]{DejaVu Sans Mono}%
}{}

% --- En-têtes et pieds de page ------------------------------------------------
\usepackage{scrlayer-scrpage}
\usepackage{lastpage}
\pagestyle{scrheadings}
\clearpairofpagestyles
\makeatletter
\ihead{\small\sffamily\@title}
\makeatother
\ohead{\small\sffamily Lycée -- SNT \& NSI}
\cfoot{\small\sffamily Page~\thepage{}~/~\pageref*{LastPage}}
\setkomafont{pagehead}{\normalfont\color{gray!70!black}}
\setkomafont{pagefoot}{\normalfont\color{gray!70!black}}

% --- Titres : couleur sobre et bon contraste ----------------------------------
\usepackage{xcolor}
\definecolor{titre}{HTML}{1F4E79}
\addtokomafont{disposition}{\color{titre}}
\addtokomafont{title}{\color{titre}}

% --- Paragraphes aérés (meilleure lisibilité) --------------------------------
\KOMAoptions{parskip=half}
\raggedright           % pas de justification : évite les espaces irréguliers

% --- Mention au-dessus des blocs Python interactifs ---------------------------
\newcommand{\pyinteractif}[1]{%
  \par\smallskip\noindent
  {\small\sffamily\color{titre}\fbox{\textbf{>\_}}~\textit{#1 -- à tester
  dans la version en ligne du cours}}\par\nopagebreak[4]%
}
